%!TEX program = xelatex

\documentclass{article}
\usepackage[UTF8]{ctex}
\usepackage{indentfirst}
\usepackage{lmodern}% http://ctan.org/pkg/lm
\usepackage{setspace}
\usepackage{verbatim}
\usepackage{amsmath}
\usepackage{amsthm}
\usepackage{graphicx}
\usepackage{subfigure}
\usepackage{booktabs}
\usepackage{tabularx}
\usepackage{multirow}
\usepackage{multicol}
\usepackage{listings}
\usepackage{relsize}
\usepackage[usenames]{color}   
\usepackage{fontspec}

\usepackage[table]{xcolor}
\usepackage{array}
\usepackage[rightcaption]{sidecap}
\usepackage{hyperref}
\usepackage{float}
\usepackage{url}

% 自己画图
\usepackage{tikz}
\usepackage{mathpazo}
\usepackage{pgfplots}

% 插入动画
\usepackage{animate}
\usepackage{caption}
\usepackage{subfigure}
\usepackage{graphicx}

% 解决目录红框问题
\hypersetup{
colorlinks=true,
linkcolor=black
}

% 代码格式设置
\lstset{
         language = python, numbers=left, 
         numberstyle=\tiny,keywordstyle=\color{blue!70},
         commentstyle=\color{red!50!green!50!blue!50},frame=shadowbox,
         rulesepcolor=\color{red!20!green!20!blue!20},basicstyle=\ttfamily
}




\newcommand{\myd}{\;\mathrm{d}} %自定义积分变量命令
\renewcommand{\listtablename}{Tables} %改变listoftable的名字
\renewcommand{\arraystretch}{1.8} %行间距相对原始间距的倍数
% \setmonofont{Consolas}
\renewcommand\refname{参考网站}
% \setlength{\tabcolsep}{18pt} % 文本到单元格左右边框之间的距离
% \setlength{\arrayrulewidth}{0.4mm} % 设置表格线宽度


\setlength{\parindent}{2em} %控制首行缩进
\addtolength{\parskip}{3pt} %\parskip 控制段落(paragraph)距离
%\singlespacing % 设置单倍行距
%\doublespacing % 设置双倍行距
%\setstretch{1.25} % 设置行距为1.25
\onehalfspacing % 1.5倍行距
\graphicspath{{./figures/}} % 指定图片所在文件夹
%\newtheorem{definition}[定义]{section}
%定制环境: {环境名}[编号延续]{显示名}[编号层次]

%\begin{doublespacing} 特定行距环境 \end{doublespacing}


\begin{document}
\begin{titlepage}
% titlepage格式可以自己调
        \vspace*{-3cm}
	
	\begin{figure}[h]
		\centering
		\includegraphics[width=0.7\linewidth]{zjdx}
	\end{figure}

	\vspace*{0.5cm}
	\begin{figure}[h]
		\centering
		\includegraphics[width=0.5\linewidth]{QSY}
	\end{figure}
	\vspace{-0.5cm}
	\begin{center}
		\Huge{\textbf{应用光学实验}}\\
		
		\Huge{\textbf{实验报告}}
	\end{center}
	
	\vspace*{0.5cm}


	\vspace*{1cm}
    \begin{center}
            \Large 实验名称\ \ \underline{\makebox[250pt]{title here}} \\ 
            \vspace{0.3cm}
            \Large 姓名 \ \ \underline{\makebox[220pt]{Peter_H}} \\ 
            \vspace{0.3cm}
            学  号\ \ \underline{\makebox[220pt]{3180xxxxxx}}\\
            \vspace{0.3cm}
            实验地点\ \ \underline{\makebox[250pt]{你猜}}\\
            \vspace{0.3cm}
            实验日期\ \ \underline{\makebox[250pt]{2020年 mm 月 dd 日}}\\
            \vspace{0.3cm}
            指导老师\ \ \underline{\makebox[250pt]{Peter_H’s Instructors}}\\
            

             
    \end{center}
        
    
\end{titlepage}

\newpage
% \setcounter{tocdepth}{5}
\tableofcontents
\thispagestyle{empty} %清楚图片格式（取消掉编号）
\listoffigures%生成插图
\listoftables%生成表格
\thispagestyle{empty} 

\newpage
\pagenumbering{arabic}%重新开始标号,阿拉伯数字形式


\section{示例}
\subsection{图片示例}
\begin{figure}[H]
\centering
\includegraphics[width=0.8\textwidth]{zjdx}
\caption{浙江大学}
\end{figure}


\subsection{表格示例}
\begin{table}[H]
  \centering
  \caption{数据记录表}
    \begin{tabular}{|c|c|c|}
    \hline
    \multicolumn{3}{|c|}{合并单元格}\\
    \hline
    \multicolumn{3}{|c|}{text}\\
    \hline
    手机焦距 $f_2'$/mm & 物到手机镜头的距离/mm &眼镜到手机的距离(透镜间隔 $d$/mm）\\
    \hline
    4.00  & 150.00  & 40.00  \\
    \hline
    像的放大倍率 $\beta$ &中间像到眼镜的位置/mm &中间像的大小$y'$/mm\\
    \hline
    0.024  & -63.25 & 5.75 \\
    \hline
    \end{tabular}
\end{table}

\subsection{代码示例}
\begin{lstlisting}
code here
set your language at line 48

\end{lstlisting}

\subsection{其他示例}
\begin{enumerate}
 \item  栗子1\cite{apple}
 \item  栗子2
 
\end{enumerate}




% \newpage
% 参考网站
\begin{thebibliography}{9} %参考文献的数量
\bibitem{apple}
  \url{https://www.apple.com.cn/}


\end{thebibliography}


\end{document}

